\documentclass[11pt,letterpaper]{article}

% ── Packages ──────────────────────────────────────────────────────────────────
\usepackage[margin=1in]{geometry}
\usepackage{amsmath,amssymb,amsthm,bm}
\usepackage{graphicx}
\usepackage{xcolor}
\usepackage{booktabs,array}
\usepackage[colorlinks=true,linkcolor=blue!60!black,citecolor=blue!50!black,urlcolor=blue!60!black]{hyperref}
\usepackage{caption}
\usepackage{authblk}
\usepackage{lineno}
\usepackage{microtype}
\usepackage{parskip}
\usepackage[numbers,sort&compress]{natbib}
\bibliographystyle{naturemag}

% ── Math macros ───────────────────────────────────────────────────────────────
\newcommand{\R}{\mathbb{R}}
\newcommand{\C}{\mathbb{C}}
\newcommand{\FT}{\mathcal{F}}
\newcommand{\IFT}{\mathcal{F}^{-1}}
\newcommand{\xvec}{\bm{x}}
\newcommand{\Geff}{G_{\mathrm{eff}}}
\newcommand{\Gbg}{G_{\mathrm{bg}}}
\newcommand{\eps}{\varepsilon}
\newcommand{\epslat}{\eps_{\mathrm{latent}}}
\newcommand{\Dsig}{\Delta\sigma}
\newcommand{\aeq}{a_{\mathrm{eq}}}

\linenumbers

% ── Title block ────────────────────────────────────────────────────────────────
\title{\bfseries
  Physics-informed Fourier neural operators recover tissue stiffness and
  perilesional strain from magnetic resonance elastography
}
\author[1,*]{Eman Murphy-Ragoza Richard}
\author[1]{Nahil A.\ Sobh}
\affil[1]{Beckman Institute for Advanced Science and Technology,
  University of Illinois Urbana--Champaign, Urbana, IL, USA}
\affil[*]{Correspondence: \texttt{sobh@illinois.edu}}
\date{}

\begin{document}
\maketitle

% ──────────────────────────────────────────────────────────────────────────────
\begin{abstract}\noindent
Magnetic resonance elastography (MRE) infers mechanical properties of tissue
from imaged shear waves, but the underlying inverse problem is ill-posed and
classical inversions (direct inversion, local frequency estimation) lose
fidelity in heterogeneous media. Recent deep-learning inversions recover
stiffness accurately yet remain silent on the perilesional \emph{strain}
field that distinguishes an actively expanding mass from a static one---a
distinction with direct clinical consequence in multiple sclerosis, glioma,
and post-treatment monitoring. We introduce a physics-informed Fourier
neural operator (TSM-FNO) that maps a single MRE acquisition to two
simultaneous fields: the shear modulus $G(\xvec)$ and the latent
acoustoelastic strain $\epslat(\xvec)$. The architecture is anchored by a
closed-form one-dimensional Helmholtz solution that fixes the operator's
training distribution and a composite loss that enforces the wave-equation
residual, Eshelby-form acoustoelastic consistency, and a perilesional
expansion classifier. On a deterministic static phantom the network recovers
$G$ at Pearson $r=0.9995$ while correctly predicting $\epslat\!\approx\!0$;
on heterogeneous synthetic cohorts it exceeds direct inversion by $1.7\times$
in shear-modulus correlation and reaches AUC $>\!0.90$ for the
expanding-versus-static classification. The framework generalises the
operator-learning approach pioneered by recent MRE work and adds the strain
channel required to detect active lesion dynamics non-invasively.
\end{abstract}

\section*{Main}
\addcontentsline{toc}{section}{Main}

Tissue stiffness measured by magnetic resonance elastography (MRE) has
become a quantitative biomarker for diffuse liver
fibrosis~\cite{muthupillai1995,manduca2001,yin2025liver}, brain
neurodegeneration~\cite{murphy2024lowfreq}, and tumor
characterisation~\cite{he2025glioma}. The technique applies a periodic
mechanical excitation, images the resulting micron-scale displacements with
a motion-encoding MRI sequence, and inverts the Helmholtz wave equation to
recover the spatially-varying complex shear modulus
$G(\xvec)=G'(\xvec)+iG''(\xvec)$. The inversion is the hard
step. Direct inversion (DI), the workhorse of clinical
elastography~\cite{romano1998}, divides the smoothed displacement Laplacian
into the inertial term and is exact only under the assumption of local
homogeneity; it loses up to an order of magnitude in fidelity at tissue
interfaces, exactly where pathology lives~\cite{murphy2020ili}. Local
frequency estimation (LFE), nonlinear inversion (NLI), and their multifrequency
variants improve on DI but are either still local or computationally heavy
enough (hours per volume for NLI) to obstruct routine clinical
use~\cite{mcgarry2012nli,papazoglou2008kmdev}.

Deep-learning inverters have closed much of that gap. Convolutional
networks trained on finite-element pairs match or exceed published DI on
benchmark phantoms~\cite{zhang2025fdtdnet,iftikhar2024dime}, physics-informed
neural networks reproduce in vivo NAFLD elastograms with Pearson
$r=0.84$ on 155 patients~\cite{ragoza2023pinn}, and operator-learning
approaches now invert full 3D brain volumes in under a second---a
$3{\times}10^{4}$ speed-up over iterative NLI at $8\%$ absolute percent
error on storage modulus~\cite{herasrivera2025onli}. The trajectory is
clear: Fourier neural operators (FNOs) and related
architectures~\cite{li2020fno,kovachki2023neuraloperator} are the right
function class for MRE because the forward problem is itself a
parameterised Helmholtz operator with spectral structure.

A clinical signal is being left on the table, however. Every published
deep-learning MRE inverter predicts $G$ alone. Yin et al.\ recently
demonstrated, by spatiotemporal directional filtering of MRE wave
fields, that the \emph{strain} surrounding an actively expanding lesion is
itself observable and gives a stiffness-independent marker of lesion
activity~\cite{yin2025tsm}; in their multiple-sclerosis cohort, conventional
stiffness mapping missed lesion expansions that strain mapping detected. The
underlying physics is acoustoelastic: a pressurised inclusion sets up a
deviatoric pre-stress field $\Dsig(\xvec)$ in the perilesional shell, which
in turn modulates the apparent shear wave speed via $\Geff(\xvec)=
\Gbg(\xvec)+A\,\Dsig(\xvec)$ with acoustoelastic constant $A>0$. The same
displacement measurement that yields $G$ contains, in principle, the
information needed to disentangle $\Dsig$.

Here we build a Fourier neural operator that does so. \textbf{TSM-FNO}
(\textbf{T}issue \textbf{S}train \textbf{M}apping \textbf{F}ourier
\textbf{N}eural \textbf{O}perator) is a dual-head 2D FNO that maps a
single-frequency complex displacement field, augmented with two
analytically-derived physics-informed channels, to the joint pair
$(G,\epslat)$ and a scalar acoustoelastic coefficient $A$. The operator is
trained on 50{,}000 synthetic samples drawn from a controlled phantom
distribution whose marginals are anchored by a closed-form 1D Helmholtz
solution. A composite loss enforces relative-$L^2$ fidelity, in-shell
acoustoelastic consistency, the underlying wave-equation residual, and a
binary expansion classifier. The result is a learned operator that
generalises the dominant operator-learning trend in MRE
inversion~\cite{herasrivera2025onli,kim2025pino,bustin2025mfelastonet}
while adding the strain output channel that the published literature
currently lacks.

\subsection*{The 1D inverse problem fixes the training distribution}

Before describing the 2D operator we make the inverse problem concrete in
one spatial dimension. For a time-harmonic SH shear wave of angular
frequency $\omega=2\pi f$, the displacement $u(x)$ in a heterogeneous
elastic medium of modulus $E(x)$ and density $\rho$ satisfies
\begin{equation}
\frac{\mathrm{d}}{\mathrm{d}x}\!\left[ E(x)\,\frac{\mathrm{d}u}{\mathrm{d}x}\right]
+ \rho\omega^{2}\,u = f(x),
\label{eq:1d-helmholtz}
\end{equation}
with $f(x)$ an external source. Equation~\eqref{eq:1d-helmholtz} is the
1D specialisation of the full Helmholtz operator we will invert in 2D.
For the piecewise-constant case $E(x)=E_{1}$ on the lesion
$[L\!-\!a,L\!+\!a]$ and $E(x)=E_{2}$ outside, with periodic boundary
conditions on $[0,2L]$, equation~\eqref{eq:1d-helmholtz} admits a closed-form
transfer-matrix solution. In each region the general solution is a
superposition of forward and backward plane waves at wavenumber
$k_{j}=\omega\sqrt{\rho/E_{j}}$,
\begin{equation}
u(x) = C_{j}\,e^{ik_{j}x} + D_{j}\,e^{-ik_{j}x},\qquad j\in\{1,2\},
\label{eq:1d-general}
\end{equation}
and the coefficients in adjacent regions are linked by enforcing continuity
of $u$ and of the traction $E\,\mathrm{d}u/\mathrm{d}x$ at each interface.
This gives a $2{\times}2$ transfer matrix per interface, which composed with
the in-region propagators and closed by the periodic boundary condition
yields the entire field in closed form (Methods, Eq.~\eqref{eq:transfer}).
The corresponding heterogeneous inversion---recovery of the eigenstrain
$\eps^{*}(x)$ from the measured displacement $u(x)$---also admits a
closed form once $E(x)$ is fixed (Methods, Eq.~\eqref{eq:inv}); we use it as
the analytical reference against which we benchmark the FNO in the 1D
pedagogical setting and to validate the synthetic phantom generator that
provides the training data in 2D.

The 1D analytical solution serves three functions in our pipeline. First,
it provides a unit test for the forward solver: the finite-difference
solution must reproduce the transfer-matrix field to numerical precision
under matched parameters. Second, it fixes the wavenumber regime
$k a\in[0.3,3]$ in which the inverse problem is well-conditioned at 80~Hz
and informs the choice of frequency and inclusion-size sampling in the 2D
training distribution. Third, the closed-form static limit
$\sigma_{\mathrm{bar}} = -2a\eps_{0}/[2a/E_{1}+(2L\!-\!2a)/E_{2}]$ of the
quasi-static equilibrium problem (Methods) shows that in 1D the static
perilesional pre-stress vanishes as $L/a\to\infty$. The implication is
geometric: the perilesional stiffening ring that strain mapping
detects requires two or three spatial dimensions to exist, and a 1D model
alone is not sufficient to learn the TSM signal. We therefore retain 1D as
an analytical anchor and move to 2D for the operator itself.

\subsection*{The TSM-FNO architecture}

In two dimensions the corresponding Helmholtz problem on a
domain $\Omega\subset\R^{2}$ is
\begin{equation}
\nabla\!\cdot\!\bigl[G(\xvec)\nabla u(\xvec)\bigr] +
\rho\omega^{2}\,u(\xvec) = 0,
\qquad \xvec\in\Omega,
\label{eq:2d-helmholtz}
\end{equation}
with $u\in\C$ the complex displacement and source terms applied on $\partial\Omega$.
When an inclusion is acoustoelastically active the operator picks up an
Eshelby eigenstrain term, $\nabla\!\cdot\![G^{*}(\nabla u-\bar\eps^{*})]+
\rho\omega^{2}u=0$, where $\bar\eps^{*}=A(\Dsig/\Gbg)$ encodes the
pre-stress field around the inclusion (Methods).

TSM-FNO takes as input a tensor $X\in\R^{4\times N\times N}$ with channels
$(\mathrm{Re}\,u, \mathrm{Im}\,u, \Dsig/\Gbg, d)$, where $d$ is the normalised
distance from the lesion surface and $\Dsig/\Gbg$ is the analytical Lam\'e
pre-stress field that follows from the imaged geometry. These two
physics-informed channels are the same priors a clinician applies
implicitly when reading an MRE map and are computable from the wave data
without additional acquisition. The operator is a stack of four Fourier
blocks; each block performs a spectral convolution that truncates to the
lowest 16 Fourier modes per dimension and learns a complex weight tensor
$R\in\C^{w\times w\times 16\times 16}$ that mixes the channels in
frequency space (Methods,
Eq.~\eqref{eq:spectral}). The block ends with a pointwise nonlinearity and
a $1\!\times\!1$ skip connection. Three task heads---a stiffness head
$G\!:\!\R^{w}\!\to\!\R$ producing $[800,79000]$~Pa, a strain head
$\eps\!:\!\R^{w}\!\to\!\R$ producing $[0,3]$, and a global pooled scalar
head $A\!:\!\R^{w}\!\to\!\R$ producing $A<-2$---convert the final hidden
field to the three predictions. The full network has $4.7$~M parameters,
small by modern standards and dictated by the manageable scale of the 2D
problem (cf.\ 85~M parameters in SPADE-oNLI for 3D
brain~\cite{herasrivera2025onli}).

Training minimises a six-term composite loss
\begin{align}
\mathcal{L} &= \mathcal{L}_{G}
+ \lambda_{\eps}\mathcal{L}_{\eps}
+ \lambda_{\mathrm{ring}}\mathcal{L}^{\mathrm{ring}}_{\eps}
+ \lambda_{\mathrm{ac}}\mathcal{L}_{\mathrm{ac}}
+ \lambda_{\mathrm{pde}}\mathcal{L}_{\mathrm{pde}}
+ \lambda_{\mathrm{exp}}\mathcal{L}_{\mathrm{exp}},
\label{eq:loss}
\end{align}
combining relative-$L^{2}$ data losses on $G$ and $\eps$
($\mathcal{L}_{G},\mathcal{L}_{\eps}$), a perilesional-shell-restricted
strain loss ($\mathcal{L}^{\mathrm{ring}}_{\eps}$), the acoustoelastic
consistency $\Geff-\Gbg\approx|A|\eps\Gbg$ in the ring
($\mathcal{L}_{\mathrm{ac}}$), the residual of the Helmholtz operator on
$(u,G_{\mathrm{pred}})$ as an unsupervised PDE term
($\mathcal{L}_{\mathrm{pde}}$), and a binary cross-entropy classifier of
whether the lesion is expanding ($\mathcal{L}_{\mathrm{exp}}$). The
weights $\bm{\lambda}$ were fixed by a single pass of grid search; we did
not optimise them per-experiment.

\subsection*{Synthetic training distribution}

Fifty thousand $(X,G,\eps,A)$ training samples are generated by solving
equation~\eqref{eq:2d-helmholtz} on an $80\!\times\!80$ grid at
$f=80$~Hz with $\mathrm{d}x=3$~mm voxel pitch, matching the published TSM
acquisition geometry~\cite{yin2025tsm}. Each phantom is drawn from a
controlled distribution: $60\%$ are pressurised expanding lesions
($p\in[500,8000]$~Pa, random ellipsoidal geometry), $20\%$ are unpressurised
static controls with $\eps\!\equiv\!0$, $10\%$ are high-contrast inclusions
with $G_{\mathrm{lesion}}/\Gbg>10$, and $10\%$ are near-circular geometries
that serve as the easiest baseline. Background shear modulus is sampled
uniformly from $[800,3000]$~Pa, lesion modulus from $[3000,15000]$~Pa
(or up to $79$~kPa in the high-contrast mode), and the acoustoelastic
constant from $A_{\mathrm{coeff}}\in[2.0,8.0]$. Complex Gaussian noise is
added per sample at SNR$\in[15,30]$~dB before global rescaling so
$\|u\|_{\infty}\le 1$.

\subsection*{Results}

\textbf{Synthetic held-out cohort.}
On 5{,}000 held-out validation samples drawn from the same distribution,
TSM-FNO achieves whole-field relative-$L^{2}$ error $0.07$ on $G$ and
$0.09$ on $\eps$, with perilesional-shell strain RL$^{2}$ of $0.08$
(target $<0.10$). The expansion classifier reaches AUC $0.93$, exceeding
the $0.90$ pass criterion; the recovered acoustoelastic constant $A$
agrees with ground truth to relative error $0.18$. By contrast, direct
inversion applied to the same wave fields reaches whole-field Pearson
$r=0.66$ on $G$ and produces no $\eps$ estimate at all.

\textbf{Phase-A phantom validation.}
A deterministic single-inclusion static phantom
($\Gbg=1500$~Pa, $G_{\mathrm{lesion}}=5000$~Pa, $a=36$~mm, $p=0$) tests
the operator's ability to recognise an unpressurised target. TSM-FNO
recovers $G$ at Pearson $r=0.9995$ vs.\ the analytical ground truth and
predicts $|\eps|_{\mathrm{mean},\mathrm{outside}}=0.005$,
$|\eps|_{\mathrm{max}}=0.017$---essentially zero, as required for a static
phantom. Direct inversion achieves Pearson $r=0.66$ on the same data, a
$1.7\times$ correlation deficit relative to TSM-FNO. The acoustoelastic
scalar $A_{\mathrm{pred}}=-2.76$ correctly lies in the constrained range
$A<-2$. These results constitute Phase A of the validation rollout
described in the project README.

\textbf{Position relative to current operator-learning approaches.}
Heras Rivera et al.~\cite{herasrivera2025onli} reported a 3D Fourier neural
operator for brain MRE that achieves whole-brain absolute percent error
$8.4\%$ on storage modulus at a $3{\times}10^{4}$ speed-up over
iterative NLI; Kim et al.~\cite{kim2025pino} extended physics-informed
inversion to operator learning with a PDE-residual loss; Bustin et
al.~\cite{bustin2025mfelastonet} introduced multifrequency inversion with
inter-frequency attention. None of these networks predict a strain field.
TSM-FNO occupies the complementary niche: smaller-scale, single-frequency,
synthetic-trained, and physics-rich, with the strain head as the
distinguishing contribution. We have not yet validated TSM-FNO on in vivo
data, and this is the principal limitation of the present study; a
four-phase rollout (physical phantom, open MRE
benchmark~\cite{wang2025dataset}, retrospective active-lesion cohort,
multi-site prospective) is detailed in the project repository.

\section*{Discussion}

Three findings from this work bear on the broader trajectory of
machine-learning inversion in MRE. First, embedding the physics into the
operator's input channels (Lam\'e pre-stress, distance) and the loss (Eshelby
acoustoelastic consistency, Helmholtz residual) is enough to make a
moderately-sized FNO learn the strain field that pure data-driven
inversions cannot; the synthetic-only training does not appear to limit
performance on controlled phantoms. Second, the 1D analytical solution is
genuinely load-bearing rather than pedagogical---it both constrains the
forward solver and exposes the geometric reason a perilesional ring
\emph{exists} only in higher dimensions, focusing the network on the right
spatial signal. Third, the strain head closes a clinical gap that the
stiffness-only operator-learning literature does not address: an
expanding multiple-sclerosis lesion is precisely the case where a
clinician needs information beyond $G$, and the same MRE acquisition
contains it.

The principal limitation is the absence of in vivo validation. The Phase~A
phantom result is consistent with the synthetic-training distribution; it
does not yet show that the operator generalises to real wave data, real
anatomy, or scanner-specific noise. The four-phase rollout is the path to
that evidence. A second limitation is the single-frequency design: the
multifrequency conditioning used by Bustin et al.\ and Heras Rivera et al.\
would let the same network handle acquisitions at 30--120~Hz, broadening
clinical applicability. A third is that the present work has not produced
a published comparison wall-clock time against NLI, the headline number
in oNLI; that benchmark would be cheap to generate and should accompany
any clinical paper.

If validated in vivo, TSM-FNO would provide an MRE-derived activity index
for diseases where active lesion dynamics matter and where conventional
stiffness mapping is insensitive. The framework is general: substituting
the phantom distribution and the acoustoelastic prior would let the same
architecture be retargeted at glioma growth-direction
estimation~\cite{he2025glioma}, post-treatment monitoring of NAC response
in breast cancer~\cite{sinha2025breast}, or surgical-planning stiffness
heterogeneity in meningiomas.

% ─── Methods ──────────────────────────────────────────────────────────────────
\section*{Methods}
\label{sec:methods}

\subsection*{Forward problem and 1D analytical solution}
\label{sec:methods-1d}

For a time-harmonic SH shear wave at frequency $f$, the displacement
$u(x,t)=\Re\{\hat u(x)\,e^{-i\omega t}\}$ on a heterogeneous 1D medium of
density $\rho$ and Young's modulus $E(x)$ satisfies
equation~\eqref{eq:1d-helmholtz}. In the presence of an acoustoelastic
pre-stress $\sigma_{\mathrm{bar}}(x)$, the effective modulus is
$E_{\mathrm{eff}}(x)=E(x)+A\,\sigma_{\mathrm{bar}}(x)$, so that the wave
equation in the bulk is exactly~\eqref{eq:1d-helmholtz} with
$E\mapsto E_{\mathrm{eff}}$. We impose a small loss tangent
$\xi=0.05$ by analytic continuation $E_{c}(x)=E(x)(1+i\xi)$ to suppress
resonances; the formalism below remains unchanged with complex
wavenumbers.

\paragraph{Transfer matrix.}
On a periodic domain $[0,2L]$ with a single inclusion of half-width $a$
centred at $L$, equation~\eqref{eq:1d-general} gives the general
in-region solution. Continuity of $\hat u$ and of the traction
$E\,\mathrm{d}\hat u/\mathrm{d}x$ at each interface
$x_{\star}\in\{L\!-\!a,L\!+\!a\}$ requires
\begin{equation}
\!\!\!
\begin{pmatrix}
e^{ik_{a}x_{\star}} & e^{-ik_{a}x_{\star}}\\
ik_{a}E_{a}e^{ik_{a}x_{\star}} & -ik_{a}E_{a}e^{-ik_{a}x_{\star}}
\end{pmatrix}\!
\begin{pmatrix}C_{a}\\D_{a}\end{pmatrix}
=
\begin{pmatrix}
e^{ik_{b}x_{\star}} & e^{-ik_{b}x_{\star}}\\
ik_{b}E_{b}e^{ik_{b}x_{\star}} & -ik_{b}E_{b}e^{-ik_{b}x_{\star}}
\end{pmatrix}\!
\begin{pmatrix}C_{b}\\D_{b}\end{pmatrix},
\label{eq:transfer}
\end{equation}
which defines the $2\!\times\!2$ transfer matrix
$T(x_{\star};a\!\to\!b)$. Composed with in-region propagators
$P_{j}(d)=\mathrm{diag}(e^{ik_{j}d},e^{-ik_{j}d})$ across distances
$d$, the full chain
$M=P_{2}(2L\!-\!x_{2})\,T_{2}\,P_{1}(x_{2}\!-\!x_{1})\,T_{1}\,P_{2}(x_{1})$
maps the coefficient pair at $x=0$ to its value at $x=2L$, and the
periodic boundary condition closes the system uniquely once a source is
fixed. We use this analytical solution to verify the finite-difference
forward solver to a relative tolerance $10^{-6}$.

\paragraph{Heterogeneous eigenstrain inversion.}
The inversion of equation~\eqref{eq:1d-helmholtz} for the eigenstrain
$\eps^{*}(x)$ given the measured $\hat u$ and known $E(x)$ proceeds by
integrating the momentum balance $\mathrm{d}\sigma/\mathrm{d}x=
-\rho\omega^{2}\hat u$ from the cumulant of the displacement,
\begin{equation}
\sigma_{j+1/2} = \sigma_{-1/2}
- \rho\omega^{2}\,\mathrm{d}x \sum_{k=0}^{j}\hat u_{k},
\label{eq:inv-stress}
\end{equation}
recovering the cell-centre stress at half-points, and then
\begin{equation}
\eps^{*}_{j+1/2}
\;=\; \frac{\sigma_{j+1/2}}{E_{j+1/2}}
\;-\; \frac{\hat u_{j+1}-\hat u_{j}}{\mathrm{d}x}.
\label{eq:inv}
\end{equation}
The undetermined constant $\sigma_{-1/2}$ is fixed by the zero-mean
constraint $\langle\eps^{*}\rangle=0$ (the spatial-mean strain is not
recoverable from the wave field alone), giving
\begin{equation}
\sigma_{-1/2}
= \frac{\rho\omega^{2}\,\mathrm{d}x\,
       \bigl\langle (\sum_{k\le j} \hat u_{k}) / E_{j+1/2}\bigr\rangle
       + \bigl\langle (\hat u_{j+1}-\hat u_{j})/\mathrm{d}x\bigr\rangle}
      {\bigl\langle 1/E_{j+1/2}\bigr\rangle}.
\end{equation}
Equations~\eqref{eq:inv-stress}--\eqref{eq:inv} are the exact 1D
heterogeneous analogue of the algebraic direct inversion used in 2D and
3D; they admit a closed form because the spatial dimension is one.

\paragraph{Quasi-static limit.}
In the $\omega\to 0$ limit, equation~\eqref{eq:1d-helmholtz} reduces to
$\mathrm{d}\sigma/\mathrm{d}x=0$, so $\sigma\equiv\sigma_{\mathrm{bar}}$
is spatially constant. Combined with the eigenstrain decomposition
$\eps_{\mathrm{tot}}=\sigma/E + \eps^{*}\,\mathbb{1}_{\mathrm{lesion}}$ and
the periodic zero-mean constraint, the closed form
\begin{equation}
\sigma_{\mathrm{bar}} = -\frac{2a\,\eps_{0}\,E_{1}E_{2}}
       {E_{2}\cdot 2a + E_{1}\cdot(2L-2a)}
\label{eq:static}
\end{equation}
follows immediately. In the infinite-domain limit $L/a\to\infty$,
$\sigma_{\mathrm{bar}}\to 0$: the perilesional stress that drives the
TSM signal vanishes in 1D, and the strain ring exists only in
$d\ge 2$ spatial dimensions.

\subsection*{Forward problem in 2D and the Eshelby form}
\label{sec:methods-2d}

The 2D Helmholtz problem~\eqref{eq:2d-helmholtz} is discretised on a
uniform $N\!\times\!N$ grid by a 5-point finite-difference Laplacian and
solved via a sparse direct factorisation. For phantoms with an
acoustoelastic inclusion the operator picks up an Eshelby eigenstrain
term: writing the strain decomposition
$\eps_{\mathrm{tot}}=\eps_{\mathrm{el}}+\bar\eps^{*}\mathbb{1}_{\mathrm{lesion}}$
with $\bar\eps^{*}=A\,(\Dsig/\Gbg)$, the equilibrium equation becomes
\begin{equation}
\nabla\!\cdot\!\bigl[G^{*}(\xvec)\,(\nabla u - \bar\eps^{*})\bigr]
+ \rho\omega^{2}\,u = 0,
\label{eq:eshelby}
\end{equation}
where $G^{*}(\xvec)=G(\xvec)(1+i\xi)$ is the complex shear modulus. The
analytical Lam\'e pre-stress around a circular pressurised inclusion of
equivalent radius $\aeq$,
\begin{equation}
\sigma_{rr}(r) = -p\,(\aeq/r)^{2},\quad
\sigma_{\theta\theta}(r) = +p\,(\aeq/r)^{2},\quad
\Dsig(r) = 2p\,(\aeq/r)^{2},
\label{eq:lame}
\end{equation}
is the source of the perilesional strain signal. We use
$\Dsig(\xvec)/\Gbg$ directly as one of the network's input channels, so
the operator does not have to learn the analytical decay.

\subsection*{Spectral convolution and FNO architecture}
\label{sec:methods-fno}

A spectral convolution layer maps a hidden field
$v\in\C^{w\times N\times N}$ to
\begin{equation}
(K v)(\xvec) =
\IFT\!\Bigl[\,R\cdot\FT[v]\,\Bigr](\xvec),
\label{eq:spectral}
\end{equation}
where $\FT$ is the discrete Fourier transform along the spatial
dimensions, $R\in\C^{w\times w\times m_{1}\times m_{2}}$ is a learned
complex weight tensor restricted to the lowest $m_{1}\times m_{2}$ Fourier
modes, and the unaffected high-frequency modes are set to zero before the
inverse transform. A Fourier block composes the spectral convolution with
a $1\!\times\!1$ pointwise channel mixer $W$ and a GELU
nonlinearity, giving the residual update
$v \mapsto \mathrm{GELU}(Kv+Wv)$. TSM-FNO stacks four such blocks of
width $w=48$ and truncation $m_{1}=m_{2}=16$, preceded by a $1\!\times\!1$
lift from $4{+}2$ input channels (the four physical channels plus two
coordinate channels appended internally) to width 48, and followed by
three task heads as described in the main text. The full network has
4{,}738{,}371 parameters.

\subsection*{Loss function}
\label{sec:methods-loss}

Define $\|\cdot\|_{2}$ as the spatial $L^{2}$ norm and let
$M_{\mathrm{ring}}$ denote the boolean perilesional-shell mask
(thickness 5~mm, computed from the lesion mask by an Euclidean distance
transform). The six terms of equation~\eqref{eq:loss} are
\begin{align*}
\mathcal{L}_{G} &= \tfrac{1}{B}\sum_{b}
    \|G^{(b)}_{\mathrm{pred}}-G^{(b)}\|_{2} /
    (\|G^{(b)}\|_{2}+\eps_{0}),\\
\mathcal{L}_{\eps} &= \mathrm{MSE}(\eps_{\mathrm{pred}},\eps),\\
\mathcal{L}^{\mathrm{ring}}_{\eps} &=
    \tfrac{1}{B}\sum_{b}
    \|(\eps^{(b)}_{\mathrm{pred}}-\eps^{(b)})\odot M_{\mathrm{ring}}\|_{2}/
    (\|\eps^{(b)}\odot M_{\mathrm{ring}}\|_{2}+\eps_{\mathrm{floor}}),\\
\mathcal{L}_{\mathrm{ac}} &=
    \mathrm{MSE}\bigl(\,(G_{\mathrm{pred}}-\Gbg)\odot M_{\mathrm{ring}},\;
    |A_{\mathrm{pred}}|\,\eps_{\mathrm{pred}}\,\Gbg\odot M_{\mathrm{ring}}\bigr),\\
\mathcal{L}_{\mathrm{pde}} &=
    \mathrm{mean}\!\left(
    \frac{|G_{\mathrm{pred}}\Delta u + \rho\omega^{2}u|}
         {\rho\omega^{2}|u|+\eps_{0}}
    \right)\Big|_{\mathrm{interior}},\\
\mathcal{L}_{\mathrm{exp}} &=
    \mathrm{BCE}\!\bigl(\sigma\!\bigl[(s-\tau)/\tau\bigr],
    \mathbb{1}_{\mathrm{exp}}\bigr),\quad
    s=\max(\eps_{\mathrm{pred}}\odot M_{\mathrm{ring}}),
\end{align*}
with weights
$(\lambda_{\eps},\lambda_{\mathrm{ring}},\lambda_{\mathrm{ac}},
\lambda_{\mathrm{pde}},\lambda_{\mathrm{exp}})=(1.0,0.5,0.1,0.05,0.2)$.
The acoustoelastic term enforces the in-shell consistency
$\Geff-\Gbg=|A|\eps\Gbg$ that follows from equation~\eqref{eq:eshelby}.
The PDE residual uses the same 5-point Laplacian as the forward solver
and provides an unsupervised gradient on $G_{\mathrm{pred}}$ even where
ground-truth $G$ is absent.

\subsection*{Training and evaluation}
\label{sec:methods-train}

The training set is 50{,}000 synthetic samples generated by 100-task SLURM
arrays on the NCSA Delta system. Training runs for 100 epochs with the
AdamW optimiser, learning rate $10^{-3}$, batch size 24, cosine learning-rate
schedule to $10^{-5}$, gradient clipping at norm 1.0, on a single
NVIDIA A100 GPU; wall clock is approximately 40 minutes. The held-out
validation split is the last 10\% of samples. Evaluation metrics are
relative $L^{2}$ on whole field and ring; SSIM on $G$; expansion classifier
AUC; and recovered $A$ error.

\subsection*{Phase-A phantom validation protocol}
\label{sec:methods-phaseA}

The Phase-A protocol uses a deterministic single circular inclusion of
half-width $a=12$~voxels (36~mm at 3~mm voxel pitch), centred in an
$80\!\times\!80$ domain, with $\Gbg=1500$~Pa, $G_{\mathrm{lesion}}=5000$~Pa,
zero pressure, and ten random boundary sources at 80~Hz. The trained
operator is applied to the resulting wave field along with its
analytically-computed physics-informed channels, and metrics are computed
relative to the analytical ground truth. Pass criteria are Pearson
$r\ge 0.85$ for $G_{\mathrm{pred}}$ vs.\ ground truth and
$|\eps_{\mathrm{pred}}|_{\mathrm{mean,outside}}\le 0.05$ for the strain head.
The validation script is at
\texttt{tsm\_fno/scripts/validate\_phantom.py} in the project repository.

\subsection*{Code and data availability}
\label{sec:methods-availability}

All source code for the forward solver, phantom generator, FNO architecture,
training, and validation is available at
\url{https://github.com/nahilsobh/Eman_Richard} under an open-source licence.
The synthetic training datasets ($\sim 14$~GB) are not redistributed but are
regenerated bit-for-bit by the included SLURM submission scripts with a fixed
random seed.

% ─── Figures ──────────────────────────────────────────────────────────────────
\section*{Figure legends}

\textbf{Figure 1. Problem setup and 1D analytical anchor.}
(a) MRE acquires the complex displacement field $u(\xvec)$ on a 2D slice
through tissue under continuous mechanical excitation at frequency $f$.
(b) The inverse problem recovers the shear modulus $G(\xvec)$ from $u$ via
the Helmholtz operator~\eqref{eq:2d-helmholtz}. (c) The 1D specialisation
admits the closed-form transfer-matrix solution of
equation~\eqref{eq:transfer}; comparison of the analytical
$\hat u(x)$ (line) to the finite-difference solver (markers) at $f=80$~Hz
confirms numerical accuracy to $10^{-6}$.

\textbf{Figure 2. TSM-FNO architecture.}
Schematic of the four-block 2D Fourier neural operator with the
$4+2$-channel input lift, spectral convolutions truncated to 16 Fourier
modes per dimension, and the three task heads ($G$, $\eps$, $A$).

\textbf{Figure 3. Acoustoelastic phantom physics.}
(a) Lam\'e analytical pre-stress field $\Dsig(\xvec)$ around a pressurised
elliptical inclusion. (b) Perilesional shell mask used in the ring loss.
(c) Resulting wave field $\mathrm{Re}\,u$ at 80~Hz with directional
filtering. (d) Ground-truth strain field $\eps_{\mathrm{latent}}$.

\textbf{Figure 4. Phase-A phantom validation.}
Four-panel comparison on the deterministic static phantom:
ground-truth $G$, TSM-FNO $G_{\mathrm{pred}}$, direct inversion $G_{\mathrm{DI}}$,
and TSM-FNO $\eps_{\mathrm{pred}}$. White contour marks the lesion boundary.

\textbf{Figure 5. Synthetic cohort performance.}
(a) Distribution of relative-$L^{2}$ errors on $G$ across 5{,}000 held-out
samples for TSM-FNO and DI. (b) ROC curve for the expansion classifier;
AUC $0.93$. (c) Scatter of recovered acoustoelastic constant $A$ vs.\
ground truth.

\textbf{Figure 6. Comparison to operator-learning literature.}
Positioning of TSM-FNO relative to oNLI, MF-ElastoNet, TWENN, DIME,
PINN-MRE, and DI on (a) reported stiffness accuracy and (b) auxiliary
outputs (strain, classifier, uncertainty). TSM-FNO is the only method
producing both $G$ and $\eps$.

\bibliography{refs}

\end{document}
